% ============================================================
%  Chương VI – Các hệ thống PKI mã nguồn mở
% ============================================================
\chapter{Các hệ thống PKI mã nguồn mở}
\label{chap:opensource}

PKI trong triển khai thực tế không chỉ là một chuẩn chứng chỉ. Nó là
một hệ sinh thái công cụ. Mỗi công cụ giải quyết một phần khác nhau của
bài toán: thao tác mật mã, quản lý CA cục bộ, tự động hóa TLS công cộng,
phát hành chứng chỉ nội bộ, hoặc tích hợp với hệ thống định danh.

Chương này khảo sát các hệ thống và công cụ mã nguồn mở theo vai trò
công nghệ. Trọng tâm là vị trí của từng công cụ trong hệ sinh thái PKI,
không phải quy trình demo hoặc thiết kế triển khai của nhóm.

% ----------------------------------------------------------
\section{Tiêu chí phân loại}
\label{sec:6.1}

Một công cụ PKI có thể được phân loại theo vai trò chính. Nhóm thứ nhất
là \textit{cryptographic toolkit}, gồm thư viện và CLI để tạo khóa, tạo
CSR, ký, chuyển đổi và kiểm tra chứng chỉ. OpenSSL thuộc nhóm này.

Nhóm thứ hai là công cụ quản lý chứng chỉ cục bộ. Chúng phù hợp với lab,
học tập hoặc hệ thống nhỏ cần giao diện trực quan. XCA là ví dụ tiêu
biểu vì nó quản lý khóa, CSR, chứng chỉ, CRL và template trong cơ sở dữ
liệu cục bộ \cite{XCA}.

Nhóm thứ ba là ACME client và public CA tự động hóa. Let’s Encrypt là CA
công cộng miễn phí và tự động. Certbot là ACME client dùng để yêu cầu,
cài đặt và gia hạn chứng chỉ từ các CA tương thích ACME
\cite{LetsEncryptDocs,CertbotDocs,RFC8555}.

Nhóm thứ tư là CA platform. Đây là các hệ thống quản trị PKI có profile,
quy trình đăng ký, phát hành, thu hồi, API và tích hợp vận hành. EJBCA
và Dogtag nằm trong nhóm này \cite{EJBCA,DogtagPKI}.

Nhóm thứ năm là internal PKI hoặc secret management. Vault PKI Secrets
Engine phát hành chứng chỉ X.509 nội bộ qua API và chính sách truy cập,
thường dùng cho dịch vụ nội bộ hoặc mTLS \cite{VaultPKI}.

Nhóm cuối là identity và authentication stack. OAuth 2.0 và OpenID
Connect không phải nền tảng PKI. Chúng dùng TLS, chữ ký số, JWT, JWS và
JWKS trong hệ xác thực và phân quyền \cite{RFC6749,OIDCCore,RFC7515,RFC7517,RFC7519}.

Các tiêu chí so sánh gồm chức năng CA, giao diện, API, ACME, vòng đời
chứng chỉ, thu hồi, OCSP/CRL, audit, tích hợp định danh và độ phức tạp
vận hành. Một công cụ tốt trong lab chưa chắc phù hợp để vận hành CA có
chính sách. Ngược lại, nền tảng enterprise thường quá nặng cho nhu cầu
kiểm thử cục bộ.

% ----------------------------------------------------------
\section{OpenSSL}
\label{sec:6.2}

OpenSSL là thư viện và dòng lệnh nền tảng cho nhiều thao tác mật mã.
Trong bối cảnh PKI, hai lệnh quan trọng là \texttt{openssl x509} và
\texttt{openssl verify}. Tài liệu OpenSSL mô tả \texttt{x509} là lệnh
hiển thị và ký chứng chỉ. Nó có thể xem trường chứng chỉ, chuyển đổi PEM
và DER, tạo chứng chỉ từ PKCS~\#10 request, tự ký chứng chỉ và ký bằng
CA key \cite{OpenSSLX509}.

Lệnh \texttt{verify} kiểm tra chuỗi chứng chỉ. Nó nhận trust anchor qua
\texttt{-CAfile}, \texttt{-CApath} hoặc nguồn tin cậy tương ứng. Nó cũng
có tùy chọn kiểm tra hostname, email, IP và mục đích sử dụng chứng chỉ
\cite{OpenSSLVerify}. Vì vậy, OpenSSL là công cụ cơ sở để quan sát và
kiểm tra hành vi PKI.

Một số lệnh ngắn cho thấy vai trò công nghệ của OpenSSL:

\begin{lstlisting}[language=bash, caption={Một số thao tác PKI với OpenSSL}, label={lst:openssl-pki}]
openssl x509 -in server.crt -text -noout
openssl verify -CAfile ca-chain.pem server.crt
openssl x509 -req -in server.csr -CA ca.crt -CAkey ca.key -out server.crt
\end{lstlisting}

Giới hạn của OpenSSL nằm ở tầng quản trị. Nó không tự cung cấp portal RA,
profile người dùng, phê duyệt, audit tập trung, HSM policy hoặc API vòng
đời chứng chỉ. Có thể xây dựng CA đơn giản bằng OpenSSL, nhưng phần quản
trị sẽ nằm trong script và quy trình bên ngoài.

% ----------------------------------------------------------
\section{XCA}
\label{sec:6.3}

XCA là công cụ giao diện đồ họa để xây dựng và quản trị các đối tượng
PKI dựa trên X.509. Trang chính thức mô tả XCA quản lý chứng chỉ, CSR,
khóa riêng RSA/DSA/EC, smartcard và CRL. Dữ liệu được lưu trong cơ sở dữ
liệu như SQLite, MySQL, PostgreSQL hoặc SQL Server qua ODBC
\cite{XCA}.

XCA phù hợp với môi trường local và lab. Người dùng có thể tạo CA nhỏ,
tạo sub-CA, dùng template cho chứng chỉ và xem chuỗi chứng chỉ trong giao
diện. Điều này làm XCA dễ tiếp cận hơn OpenSSL khi mục tiêu là học cấu
trúc PKI hoặc quản lý một kho chứng chỉ nhỏ.

So với OpenSSL, XCA mạnh ở trực quan hóa và quản lý local certificate
store. OpenSSL mạnh hơn ở CLI, scripting và kiểm tra tự động trong
pipeline. Cả hai đều không nên được mô tả như CA platform enterprise.
Chúng thiếu các lớp vận hành như RA workflow, phân quyền quản trị, API
chuẩn hóa và audit tập trung.

% ----------------------------------------------------------
\section{EJBCA}
\label{sec:6.4}

EJBCA là trọng tâm của nhóm CA platform mã nguồn mở. Tài liệu chính thức
mô tả EJBCA như hệ thống có hướng dẫn cho CA, RA, vận hành, tích hợp và
khắc phục sự cố \cite{EJBCA}. So với OpenSSL hoặc XCA, EJBCA không chỉ
ký chứng chỉ. Nó quản lý chính sách phát hành và vòng đời chứng chỉ.

Trong EJBCA, certificate profile là template giới hạn cấu trúc chứng chỉ.
Profile có thể quy định thuật toán khóa, kích thước khóa, mục đích chứng
chỉ, Key Usage, Extended Key Usage, Certificate Transparency và CA được
phép phát hành \cite{EJBCAProfiles}. End entity profile là template cho
subject DN và thông tin của chủ thể. Nó cũng giới hạn certificate profile
mà một end entity được dùng khi đăng ký \cite{EJBCAEndEntityProfiles}.

Cách tách certificate profile khỏi end entity profile cho phép EJBCA mô
hình hóa chính sách CA rõ hơn công cụ local. Certificate profile mô tả
loại chứng chỉ được phát hành. End entity profile mô tả đối tượng được
đăng ký và quyền dùng profile. RA hoặc quy trình enrollment có thể dựa
vào hai lớp này để kiểm soát yêu cầu cấp chứng chỉ.

EJBCA có REST interface cho quản trị CA và quản lý chứng chỉ. Tài liệu
REST nêu các nhóm API cho ACME, CA management, certificate, end entity,
approval, crypto token và ConfigDump. REST API mặc định bị tắt và có thể
dùng xác thực bằng client certificate hoặc OAuth \cite{EJBCAREST}.

EJBCA cũng hỗ trợ ACME. Tài liệu ACME của EJBCA ghi nhận việc triển khai
RFC~8555, hỗ trợ order, authorization, challenge, finalize, revocation,
key change, tải chứng chỉ và renewal information. Các challenge được nêu
gồm \texttt{http-01}, \texttt{dns-01}, \texttt{tls-alpn-01} và một số
kiểu mở rộng \cite{EJBCAACME}.

Về thu hồi, EJBCA cung cấp OCSP responder. Tài liệu OCSP mô tả endpoint
\texttt{/ejbca/publicweb/status/ocsp}, OCSP Key Bindings, cấu hình AIA
trong certificate profile và transaction logs cho request/response
\cite{EJBCAOCSP}. Những chức năng này đặt EJBCA gần với môi trường CA
vận hành thực tế hơn công cụ chỉ tạo chứng chỉ.

Tài liệu EJBCA cũng có phần tích hợp Hardware Security Modules trong mục
integration \cite{EJBCA}. Điều này quan trọng vì private key của CA là
tài sản có rủi ro cao. Trong hệ thống nghiêm túc, CA key thường cần cơ
chế bảo vệ và chính sách vận hành riêng.

Vị trí công nghệ của EJBCA là CA platform hoàn chỉnh. Nó phù hợp khi cần
mô phỏng hoặc vận hành CA có policy, profile, enrollment, thu hồi, OCSP,
API và tích hợp bảo mật. Đổi lại, độ phức tạp cao hơn OpenSSL và XCA.
Người vận hành phải hiểu cấu hình CA, profile, quyền quản trị, lưu khóa
và quy trình thu hồi.

% ----------------------------------------------------------
\section{Let’s Encrypt và Certbot}
\label{sec:6.5}

Let’s Encrypt là CA công cộng miễn phí, tự động và mở do Internet
Security Research Group vận hành \cite{LetsEncryptDocs}. Nó thuộc Web
PKI công cộng, không phải CA nội bộ cho tên riêng hoặc định danh doanh
nghiệp. Certbot là ACME client phổ biến để lấy, cài đặt, gia hạn, thu
hồi, xóa và liệt kê chứng chỉ \cite{CertbotDocs}.

ACME định nghĩa cách client tự động làm việc với CA. RFC~8555 mô tả luồng
điển hình: tạo hoặc dùng account, tạo order, lấy authorization, hoàn
thành challenge, finalize order bằng CSR và tải chứng chỉ đã cấp
\cite{RFC8555}. Vì vậy, ACME chuyển nhiều thao tác cấp TLS certificate
thành giao thức tự động.

Các challenge ACME chứng minh quyền kiểm soát identifier. HTTP-01 dùng
file dưới đường dẫn \texttt{/.well-known/acme-challenge/} và chỉ chạy
trên cổng 80. DNS-01 dùng TXT record dưới
\texttt{\_acme-challenge}. TLS-ALPN-01 chạy qua TLS trên cổng 443 và dùng
ALPN riêng \cite{LetsEncryptChallenges}. DNS-01 hỗ trợ wildcard
certificate, còn HTTP-01 và TLS-ALPN-01 không hỗ trợ wildcard theo tài
liệu Let’s Encrypt \cite{LetsEncryptChallenges}.

Certbot cung cấp các plugin xác thực và cài đặt. Apache và Nginx plugin
có thể lấy và cài chứng chỉ. Webroot plugin đặt challenge file trong web
root hiện có. Standalone plugin chạy web server tạm thời và cần truy cập
cổng 80. DNS plugin dùng TXT record và cần thiết cho wildcard certificate
\cite{CertbotDocs}.

Gia hạn là điểm mạnh vận hành của Certbot. Lệnh \texttt{certbot renew}
kiểm tra các chứng chỉ đã biết và gia hạn chứng chỉ gần hết hạn. Hầu hết
cài đặt Certbot có scheduled task chạy \texttt{certbot renew}. Certbot
cũng hỗ trợ hook trước, sau và sau khi triển khai để dừng, khởi động hoặc
nạp lại dịch vụ \cite{CertbotDocs}.

Giới hạn chính của mô hình này là nó phụ thuộc vào quyền kiểm soát tên
miền. Let’s Encrypt cũng áp dụng rate limit. Tài liệu hiện hành nêu giới
hạn 50 chứng chỉ mới cho mỗi registered domain trong 7 ngày và 5 chứng
chỉ cho cùng exact identifier set trong 7 ngày \cite{LetsEncryptRateLimits}.
Do đó, tự động hóa cần môi trường staging và cấu hình gia hạn cẩn thận.

% ----------------------------------------------------------
\section{Vault PKI}
\label{sec:6.6}

Vault PKI Secrets Engine thuộc nhóm internal PKI và secret management.
Tài liệu HashiCorp mô tả engine này sinh chứng chỉ X.509 động. Dịch vụ có
thể lấy chứng chỉ qua cơ chế xác thực và phân quyền của Vault, thay vì
luồng CSR thủ công với CA ngoài \cite{VaultPKI}.

Vault PKI thường dùng role để ràng buộc chính sách phát hành. Role có thể
giới hạn domain được phép, cho phép subdomain và đặt TTL tối đa. Chứng
chỉ được phát hành bằng cách gọi endpoint \texttt{issue} theo tên role,
và Vault trả về certificate, issuing CA, private key, key type và serial
number \cite{VaultPKIQuickStart}.

Điểm mạnh của Vault PKI là API-driven issuance và TTL ngắn. Tài liệu
HashiCorp nhấn mạnh certificate TTL ngắn giúp giảm phụ thuộc vào thu hồi
và làm CRL nhỏ hơn trong môi trường nhiều workload \cite{VaultPKI}. Mỗi
instance dịch vụ có thể nhận chứng chỉ riêng, giúp giảm rủi ro khi phải
xoay vòng một chứng chỉ dùng chung.

Vault PKI phù hợp với internal TLS, mTLS giữa service và hạ tầng ứng dụng.
Nó không thay thế public CA vì chứng chỉ do Vault phát hành chỉ được tin
bởi các hệ thống đã cài trust anchor tương ứng. Nó cũng không phải CA
platform enterprise đầy đủ theo nghĩa EJBCA hoặc Dogtag. Vai trò chính
của Vault là phát hành bí mật và chứng chỉ ngắn hạn theo policy.

% ----------------------------------------------------------
\section{OAuth/OIDC và PKI}
\label{sec:6.7}

OAuth 2.0 là authorization framework cho dịch vụ HTTP. RFC~6749 định
nghĩa các vai trò như resource owner, resource server, client và
authorization server. Client nhận access token có scope và thời hạn, thay
vì dùng trực tiếp mật khẩu của resource owner \cite{RFC6749}.

OpenID Connect bổ sung identity layer trên OAuth 2.0. OpenID Connect Core
định nghĩa ID Token là JWT chứa claim về sự kiện xác thực người dùng. ID
Token phải được ký bằng JWS và có thể được mã hóa thêm bằng JWE
\cite{OIDCCore}.

JWT là định dạng compact để mang claim giữa hai bên. RFC~7519 mô tả JWT
Claims Set và các claim như \texttt{iss}, \texttt{sub}, \texttt{aud},
\texttt{exp} và \texttt{iat}. Nội dung JWT có thể được bảo vệ bằng JWS
hoặc JWE \cite{RFC7519}. JWS cung cấp chữ ký số hoặc MAC trên payload
\cite{RFC7515}.

JWKS là tập các JSON Web Key. RFC~7517 định nghĩa JWK Set là JSON object
có mảng \texttt{keys}. Trường \texttt{kid} giúp bên xác minh chọn đúng
khóa khi có nhiều khóa hoặc khi xoay khóa \cite{RFC7517}. Vì vậy,
resource server có thể lấy public key từ JWKS để kiểm tra chữ ký token.

Ranh giới với PKI cần được giữ rõ. PKI quản lý ràng buộc public key với
identity qua certificate và chuỗi CA. OAuth/OIDC quản lý authorization,
identity token và xác minh token. Điểm giao là TLS certificate bảo vệ kênh
đến IdP hoặc API, IdP ký token bằng private key, và service xác minh token
bằng public key.

Một điểm giao trực tiếp hơn là mutual TLS trong OAuth. RFC~8705 định
nghĩa \texttt{tls\_client\_auth} cho client authentication dựa trên
X.509 mutual TLS và certificate-bound access token. Với token dạng này,
resource server phải so khớp certificate trình bày qua mTLS với binding
trong token \cite{RFC8705}.

% ----------------------------------------------------------
\section{Dogtag và FreeIPA}
\label{sec:6.8}

Dogtag Certificate System là enterprise-class open source Certificate
Authority. Tài liệu Dogtag mô tả nó là full PKI platform quản lý vòng đời
chứng chỉ, gồm issuance, revocation, retrieval và profiles. Nó cũng hỗ
trợ CRL generation, OCSP responder, SCEP và local Registration Authority
\cite{DogtagPKI}.

Dogtag phù hợp với môi trường cần CA nội bộ có chức năng enterprise.
So với EJBCA, Dogtag thường được nhìn trong hệ sinh thái Linux doanh
nghiệp và Red Hat. Nó mạnh hơn OpenSSL/XCA vì có vòng đời chứng chỉ,
profile, thu hồi và dịch vụ PKI. Đổi lại, vận hành Dogtag đòi hỏi hiểu
nhiều thành phần hệ thống hơn.

FreeIPA là hệ thống identity management cho môi trường Linux/UNIX. Tài
liệu FreeIPA mô tả nó tích hợp identity, authentication, authorization và
account management. Nó quản lý user, group, host và các đối tượng bảo mật
mạng. FreeIPA dùng Dogtag Certificate System để cung cấp CA và RA
functions \cite{FreeIPA}.

Vị trí của FreeIPA khác Dogtag thuần PKI. FreeIPA đặt certificate service
trong hệ định danh gồm Kerberos, LDAP, DNS, web UI và CLI \cite{FreeIPA}.
Nó phù hợp khi tổ chức muốn quản lý user, host, service và certificate
trong cùng một miền nội bộ. Mức phức tạp cao hơn công cụ local, nhưng đổi
lại có tích hợp định danh tập trung.

% ----------------------------------------------------------
\section{So sánh hệ sinh thái}
\label{sec:6.9}

\begin{center}
\scriptsize
\setlength{\tabcolsep}{2pt}
\begin{longtable}{@{} >{\raggedright\arraybackslash}p{0.10\textwidth}
    >{\raggedright\arraybackslash}p{0.10\textwidth}
    >{\raggedright\arraybackslash}p{0.11\textwidth}
    >{\raggedright\arraybackslash}p{0.10\textwidth}
    >{\raggedright\arraybackslash}p{0.09\textwidth}
    >{\raggedright\arraybackslash}p{0.10\textwidth}
    >{\raggedright\arraybackslash}p{0.09\textwidth}
    >{\raggedright\arraybackslash}p{0.09\textwidth}
    >{\raggedright\arraybackslash}p{0.09\textwidth}
    >{\raggedright\arraybackslash}p{0.10\textwidth} @{} }
    \caption{So sánh một số công cụ và hệ thống PKI mã nguồn mở}
    \label{tab:opensource-pki-tools} \\
    \toprule
    \textbf{Công cụ} & \textbf{Nhóm} & \textbf{Vai trò chính}
        & \textbf{Môi trường} & \textbf{CA} & \textbf{API}
        & \textbf{Thu hồi} & \textbf{Định danh}
        & \textbf{Độ phức tạp} & \textbf{Giới hạn chính} \\
    \midrule
    \endfirsthead
    \toprule
    \textbf{Công cụ} & \textbf{Nhóm} & \textbf{Vai trò chính}
        & \textbf{Môi trường} & \textbf{CA} & \textbf{API}
        & \textbf{Thu hồi} & \textbf{Định danh}
        & \textbf{Độ phức tạp} & \textbf{Giới hạn chính} \\
    \midrule
    \endhead
    OpenSSL
        & Toolkit
        & Thao tác khóa, CSR, chứng chỉ
        & Local, script
        & Rất cơ bản
        & CLI
        & Qua CRL khi cấu hình
        & Không
        & Thấp
        & Không có RA workflow hoặc audit tập trung. \\
    XCA
        & Local management
        & Quản lý CA và certificate store cục bộ
        & Lab, học tập
        & CA cục bộ
        & Thấp
        & Có CRL
        & Không
        & Thấp đến vừa
        & Không phải enterprise CA hoặc API-first system. \\
    EJBCA
        & CA platform
        & Quản trị CA có profile và lifecycle
        & CA nội bộ, CA có policy
        & Đầy đủ
        & REST, ACME
        & OCSP, CRL
        & RA, OAuth cho REST
        & Cao
        & Cần cấu hình profile, quyền và vận hành khóa. \\
    Let’s Encrypt/Certbot
        & Public ACME
        & Tự động hóa TLS public
        & Web PKI công cộng
        & CA công cộng
        & ACME, plugin
        & ACME revocation
        & Domain validation
        & Vừa
        & Phụ thuộc domain control và rate limit. \\
    Vault PKI
        & Internal PKI
        & Cấp chứng chỉ nội bộ ngắn hạn
        & Service nội bộ, mTLS
        & CA nội bộ
        & HTTP API, role
        & CRL, TTL ngắn
        & Vault auth/policy
        & Vừa
        & Không thay public CA hoặc CA enterprise. \\
    OAuth/OIDC stack
        & Identity/auth
        & Xác minh token và identity
        & IdP, API
        & Không phải CA
        & JWKS, token validation
        & Không quản lý chứng chỉ
        & Rất mạnh
        & Vừa đến cao
        & Quản lý token, không quản lý PKI. \\
    Dogtag/FreeIPA
        & Enterprise identity + PKI
        & CA và certificate trong identity domain
        & Linux enterprise
        & Đầy đủ trong Dogtag
        & Web, CLI
        & CRL, OCSP
        & LDAP, Kerberos
        & Cao
        & Phức tạp và gắn với hệ sinh thái Linux. \\
    \bottomrule
\end{longtable}
\end{center}

Có thể nhìn hệ sinh thái này theo tầng. OpenSSL và XCA nằm ở tầng local
PKI và kiểm thử. Certbot và Let’s Encrypt giải quyết tự động hóa TLS
trong Web PKI công cộng. EJBCA đại diện cho CA platform hoàn chỉnh, nơi
policy và vòng đời chứng chỉ là trung tâm.

Vault PKI phù hợp với chứng chỉ nội bộ ngắn hạn cho application
infrastructure. OAuth/OIDC không phải PKI platform, nhưng dùng chữ ký số,
JWKS, mTLS và TLS certificate trong xác thực hiện đại. Dogtag và FreeIPA
đặt PKI vào hệ identity enterprise, nơi user, host, service và certificate
được quản lý cùng nhau.
